\chapter{Нейрон --- не один прилад}
% full version: chapters/18_eetypes.tex

У першому розділі ми розсортували нейрони за тим, яку речовину вони викидають. Інженер розсортував би їх інакше: яким приладом працює кожен із них. І виявляється, що під однаковим виглядом ховається цілий ящик радіодеталей.

\section*{Ящик радіодеталей}

Діряве відро з другого розділу --- це накопичувач із витоком: він додає сигнали за коротке вікно. З дуже маленьким вікном він перетворюється на детектор збігів. Нейрон, який відповідає лише на зміни й одразу замовкає, --- це фільтр, що пропускає зміни. Нейрон, який працює сам, --- метроном, генератор; якщо його темп підкручує вхід, це генератор, керований напругою. Нейрон, що видає пачки клацань, --- генератор пачок, який пише «тире». Нейрон, що клацає в такт звуковій хвилі, --- детектор фази.

\pencilfig{18}{\pencilart{18}}{Під однаковим виглядом ховається цілий ящик радіодеталей: резистор, конденсатор, котушка, генератор.}

Є й прилади, про які ми раніше не говорили. \emph{Резонатор} схожий на гойдалку: він найкраще відповідає на поштовхи певного ритму, як гойдалка, що розгойдується, лише якщо штовхати її вчасно. Усередині нейрона роль пружини грає повільний струм, який опирається змінам. \emph{Накопичувач без витоку} тримає значення довго, як конденсатор без витоку, --- так працює мережа, що утримує положення очей, коли ви дивитеся вбік. Але для цього зворотний зв'язок має бути підігнаний майже ідеально, і це дороге налаштування. \emph{Засувка} --- нейрон, який короткий поштовх переводить у довгу роботу, доки його не вимкнуть, як вимикач із фіксацією; у нейронів, що керують м'язами, такий режим вмикає серотонін.

\section*{Прилад-хамелеон}

Головне відкриття цього сортування: це не різні деталі, а різні \emph{режими} однієї й тієї самої. Той самий нейрон може бути накопичувачем удень і генератором пачок уночі, залежно від того, що йому кажуть радіостанції. Інженер назвав би мозок перепрограмовуваною логічною матрицею: схема одна, а призначення деталей змінюють налаштування.

Що нейрони вміють працювати в усіх цих режимах, виміряно. Як мозок користується перемиканням режимів для обчислень --- здебільшого гіпотеза.

\fullversion{18}
